% Self-contained source. Compile twice with pdfLaTeX, or use latexmk -pdf.
% Bibliography is embedded. No external figures or data files are required.
\documentclass[11pt]{article}
\usepackage[T1]{fontenc}
\usepackage[utf8]{inputenc}
\usepackage{lmodern}
\usepackage[margin=1in]{geometry}
\usepackage{amsmath,amssymb,amsthm,mathtools}
\usepackage{microtype,booktabs,array}
\usepackage[hidelinks]{hyperref}
\hypersetup{pdftitle={Period Separation and the Single-Period Conjecture},
pdfsubject={An application of binary necklace partitions to the binary Ducci map}}
\newtheorem{theorem}{Theorem}[section]
\newtheorem{proposition}[theorem]{Proposition}
\newtheorem{lemma}[theorem]{Lemma}
\newtheorem{corollary}[theorem]{Corollary}
\newtheorem{conjecture}[theorem]{Conjecture}
\theoremstyle{definition}
\newtheorem{definition}[theorem]{Definition}
\newtheorem{example}[theorem]{Example}
\theoremstyle{remark}
\newtheorem{remark}[theorem]{Remark}
\numberwithin{equation}{section}
\newcommand{\F}{\mathbb F_2}
\newcommand{\Z}{\mathbb Z}
\newcommand{\Spec}{\mathcal S}
\DeclareMathOperator{\ord}{ord}
\DeclareMathOperator{\lcm}{lcm}
\DeclareMathOperator{\Fix}{Fix}
\setlength{\emergencystretch}{2em}
\allowdisplaybreaks[1]
\title{Period Separation and the Single-Period Conjecture\\[.35em]
\large An application of binary necklace partitions}
\author{}
\date{}

\begin{document}
\maketitle
\vspace{-2.5em}
\begin{abstract}
The single-period conjecture for the binary Ducci map asks whether an odd
width at which every nonzero periodic word has the same temporal period
must be prime. We apply the partition into binary-necklace components to
this question. If $2^s\equiv-1\pmod n$, the classical identity
$D^{2^s-1}=R^{-1}$ on the periodic subspace implies that a word of temporal
period $\tau$ has spatial period $\gcd(n,\tau)$. Distinct spatial-period
classes therefore have disjoint temporal spectra. Every divisor $d>1$ of
$n$ occurs, yielding at least $d(n)-1$ nonzero periods, where $d(n)$ is
the number of positive divisors of $n$. We give the exact cardinality of
each spatial class and of its full basin in the power set. Combining this
separation with the finite-field degree condition proves that a composite
single-period width, if one exists, must have odd multiplicative order of
$2$. A polynomial-gcd criterion then certifies period diversity and counts
both parts of the resulting partition without selecting words. At Mersenne
widths it gives a complete ordinary-period census, including exact counts
for $n=2047$. The remaining question is a uniform obstruction for composite
widths satisfying the odd-order and constant-degree conditions.
\end{abstract}

\noindent\textbf{Keywords.} Binary necklaces; Ducci map; temporal periods;
spatial periods; finite fields; overpseudoprimes; M\"obius inversion.

\section{The question and the contribution}\label{sec:introduction}

Let $n>1$ be odd. On $V_n=\F^n$, let $R$ rotate coordinates one position
to the right and put $D=I+R$. Thus $D$ replaces each bit by its sum modulo
two with its preceding neighbor. On binary entries this is the absolute
difference operation of the Ducci map. A word $v$ is \emph{periodic} if
$D^\tau v=v$ for some $\tau>0$; its \emph{ordinary temporal period} is the
least such $\tau$. Define
\begin{equation}\label{eq:spectrum}
 \Spec(n)=\{\operatorname{per}_D(v):v\ne0,\ v\text{ is periodic}\}.
\end{equation}
The zero fixed point is excluded. It is the only fixed word, since
$Dv=v$ would give $Rv=0$.

\begin{conjecture}[Single-period conjecture]\label{conj:single}
If $n>1$ is odd and $|\Spec(n)|=1$, then $n$ is prime.
\end{conjecture}

Breuer, L\"otter, and van der Merwe pose this as Question~6.1
in \cite{BLM}, counting the zero fixed point's period $1$ as well: do exactly two periods
force prime width? The conjecture is also recorded in OEIS A138004
\cite{OEIS}. The converse is false; prime width need not give a single
nonzero period. Nor does a single period mean a single cycle: many cycles
can have the same length.

The companion papers \emph{Consequences of Binary Necklaces}
\cite{Introduction} and \emph{The Partitioning of Power Sets Into Graphs}
\cite{Partition} pass from words to their rotation classes and then to
components of the induced functional graph. Here that partition becomes
a tool for a specific period question. The essential point is to retain
spatial period and rotational drift. The unlabeled graph shape alone does
not determine a word's temporal period.

We develop three consequences. First, when a power of $2$ is $-1$ modulo
$n$, a temporal period determines its spatial period. The component
partition then separates at least $d(n)-1$ temporal periods, where $d(n)$
counts the positive divisors of $n$. Second, a finite-field degree argument
extends the exclusion to \emph{every composite width with even}
$\ord_n(2)$. Third, fixed subspaces give exact polynomial certificates of
period diversity, including for widths outside the first theorem.

The polynomial description of Ducci periods is established
\cite[Theorems~3.7 and~3.9]{BLM}, and the operator identity behind our first
theorem is classical \cite[proof of Theorem~2]{Ehrlich}. We derive their
spatial-period refinement, its mass formulas, and the resulting reduction
of Conjecture~\ref{conj:single}. All required arguments are included.
These are presented as deductions from the established framework; no
claim of priority for the corollaries is made.

\section{The periodic part and the necklace partition}\label{sec:setup}

\subsection{The periodic subspace}

Identify a word $v$ with $\sum_{i=0}^{n-1}v_i x^i$ modulo $x^n+1$.
Then $R$ and $D$ are multiplication by $x$ and $1+x$. Set
\begin{equation}\label{eq:Q}
 Q_n(x)=\frac{x^n+1}{x+1}.
\end{equation}
The polynomial $x^n+1$ is squarefree because $n$ is odd. The Chinese
remainder theorem gives
\begin{equation}\label{eq:crt}
 V_n\simeq\F\oplus W_n,
 \qquad W_n=\F[x]/Q_n(x).
\end{equation}
Evolution kills the first coordinate and is invertible on $W_n$.
Consequently the periodic words correspond precisely to $\{0\}\oplus W_n$;
we identify this subspace with $W_n$. In word coordinates it is the
even-weight subspace: the first CRT coordinate is evaluation at $x=1$.
It has $2^{n-1}$ elements.

Let $\rho:V_n\to W_n$ be the periodic-coordinate projection. Each
$w\in W_n$ has two lifts under $\rho$, one periodic and one transient.
The eventual temporal period of either lift equals the period of $w$.
Thus a union of periodic orbits containing $a$ core words corresponds to
$2a$ starting words in the full power set. In particular, the zero core
has a basin of size two.

\subsection{Three different periods}

For a nonzero periodic word $w$, write
\begin{align}
 d(w)&=\min\{d>0:R^d w=w\},\label{eq:spatial}\\
 t(w)&=\min\{t>0:D^t w\in\langle R\rangle w\},\label{eq:necklace}\\
 \tau(w)&=\min\{\tau>0:D^\tau w=w\}.\label{eq:ordinary}
\end{align}
These are its spatial period, necklace period, and ordinary temporal
period. If $D^t w=R^s w$, with $0\le s<d$, then
\begin{equation}\label{eq:drift}
 \tau=\frac{td}{\gcd(d,s)}.
\end{equation}
Indeed an exact return takes a multiple $kt$ of the necklace period,
and it is an exact return precisely when $d\mid ks$.

The joint orbit $\langle R,D\rangle w$ has $dt$ periodic words.
Its rotation quotient is one necklace cycle of length $t$. Together
with its transient predecessors, it gives one component of the necklace
functional graph and represents $2dt$ words of $V_n$. The core contains
$dt/\tau=\gcd(d,s)$ ordinary cycles. These facts also follow from the
coset classification in \cite{Partition}; here they follow directly by
counting the $t$ rotation orbits of size $d$.

At odd width each cycle vertex in the necklace graph has one transient
leaf. To see uniqueness, every core word has a unique periodic predecessor
and one transient predecessor. Those two predecessors are complements.
Complementation commutes with rotation, so the transient predecessors of
one core necklace form one necklace of the same spatial period.

\subsection{Local multiplicative orders}\label{sec:local}

Factor $Q_n=\prod_{i=1}^r f_i$ into distinct monic irreducibles and put
$e_i=\deg f_i$. Then
\begin{equation}\label{eq:field-crt}
 W_n\simeq\prod_{i=1}^r\mathbb F_{2^{e_i}}.
\end{equation}
In factor $i$, let $\alpha_i$ be the image of $x$ and set
$\lambda_i=\ord(1+\alpha_i)$. Multiplication by $1+\alpha_i$ gives every
nonzero element of that coordinate period $\lambda_i$. Hence
\begin{equation}\label{eq:lcm-spectrum}
 \Spec(n)=\left\{\lcm_{i\in S}\lambda_i:
          \varnothing\ne S\subseteq\{1,\ldots,r\}\right\}.
\end{equation}
Every support occurs in \eqref{eq:field-crt}. In particular,
\begin{equation}\label{eq:local-single}
 |\Spec(n)|=1\quad\Longleftrightarrow\quad
 \lambda_1=\cdots=\lambda_r.
\end{equation}
This is the standard local-order criterion \cite{BLM,OEIS}. It will be
used for the reduction in Section~\ref{sec:odd-order}; the next theorem
does not require factoring $Q_n$.

\section{Separation by spatial period}\label{sec:separation}

Suppose throughout this section that
\begin{equation}\label{eq:minus-one}
 2^s\equiv-1\pmod n\quad\text{for some }s\ge1,
 \qquad B=2^s-1.
\end{equation}
Since $n\mid2^s+1$, we have
\begin{equation}\label{eq:coprime}
 \gcd(n,B)=1.
\end{equation}

\begin{theorem}[Spatial-period separation]\label{thm:separation}
Under \eqref{eq:minus-one}, every nonzero periodic word satisfies
\begin{equation}\label{eq:separation}
 \boxed{\quad d=\gcd(n,\tau),\qquad
 t=\gcd(B,\tau)=\frac{\tau}{d},\qquad \tau=dt.\quad}
\end{equation}
In particular, $t\mid B$, $\gcd(t,n)=1$, and words with different spatial
periods cannot have the same ordinary temporal period. Each necklace
component has exactly one ordinary cycle in its core and has mass
$2\tau=2dt$ in the power set.
\end{theorem}

\begin{proof}
Repeated squaring in characteristic two gives
\[
 D^{2^s}=I+R^{2^s}=I+R^{-1}=R^{-1}D.
\]
On the periodic subspace $D$ is invertible, so cancellation yields the
classical identity
\begin{equation}\label{eq:ehrlich}
 D^B=R^{-1}\qquad\text{on }W_n.
\end{equation}
In particular $D^{nB}=I$, so $\tau\mid nB$. By \eqref{eq:coprime},
there are unique $d_0\mid n$ and $t_0\mid B$ with $\tau=d_0t_0$.
On an ordinary $D$-cycle of length $\tau$, the permutation $D^B$ has
order $\tau/\gcd(\tau,B)=d_0$. Equation~\eqref{eq:ehrlich} identifies
this order with the spatial period $d$, so
$d=d_0=\gcd(\tau,n)$.

Moreover $R=D^{-B}$ on the core. Thus rotation stays inside an ordinary
$D$-cycle, and the joint $\langle R,D\rangle$ orbit is that same cycle.
Its quotient by rotation has $\tau/d=t_0$ vertices and is one necklace
cycle. This proves the formula for $t$, as well as the one-cycle assertion.
The two lifts of each core word give mass $2\tau$.
\end{proof}

\subsection{Exact sizes of the spatial classes}

Let $d(n)$ denote the number of positive divisors of $n$, and let $\mu$
denote the integer M\"obius function. For odd $d\ge1$, define
\begin{equation}\label{eq:Ad}
 A(d)=\sum_{e\mid d}\mu(d/e)2^{e-1}.
\end{equation}
In particular $A(1)=1$.

\begin{lemma}\label{lem:spatial-count}
For every odd $n$ and every $d\mid n$, exactly $A(d)$ periodic words
have spatial period $d$. If $d>1$, then $A(d)>0$.
\end{lemma}

\begin{proof}
A word fixed by $R^d$ repeats a block of length $d$. Because $n/d$ is
odd, the repeated word has even weight exactly when its block does.
There are $2^{d-1}$ such words. Therefore
$2^{d-1}=\sum_{e\mid d}A(e)$, and M\"obius inversion gives
\eqref{eq:Ad}. For positivity, a block of odd length $d>1$ with weight
two cannot be a proper repetition: the odd repetition number would
divide two. Its repetition to length $n$ is periodic and has spatial
period exactly $d$.
\end{proof}

For $d\mid n$, $d>1$, set
\[
 \Spec_d(n)=\{\tau(w):w\in W_n,\ d(w)=d\}.
\]
Under \eqref{eq:minus-one}, Theorem~\ref{thm:separation} gives the
disjoint union
\begin{equation}\label{eq:disjoint-spectra}
 \Spec(n)=\bigsqcup_{\substack{d\mid n\\d>1}}\Spec_d(n),
 \qquad \Spec_d(n)\subseteq
 \{dt:t\mid B,\ \gcd(t,n)=1\}.
\end{equation}
The inclusion is a necessary condition, not an assertion that every
listed $dt$ occurs. Each spatial class contains $A(d)$ core words and
its full basin contains $2A(d)$ words. These masses satisfy
\begin{equation}\label{eq:mass-accounting}
 2+\sum_{\substack{d\mid n\\d>1}}2A(d)=2^n.
\end{equation}
This equality accompanies a disjoint partition: the spatial period of
each core word is unique, and every starting word has a unique core
coordinate.

\begin{corollary}\label{cor:divisor-bound}
If \eqref{eq:minus-one} holds, then
\begin{equation}\label{eq:divisor-bound}
 |\Spec(n)|\ge d(n)-1.
\end{equation}
Consequently Conjecture~\ref{conj:single} holds for every odd width at
which a power of $2$ is $-1$ modulo $n$.
\end{corollary}

\begin{proof}
Every set in \eqref{eq:disjoint-spectra} is nonempty by
Lemma~\ref{lem:spatial-count}. If $n$ is composite, $d(n)-1\ge2$.
\end{proof}

\begin{corollary}[An infinite family]\label{cor:three-powers}
At width $n=3^k$, $k\ge1$, there are at least $k$ distinct nonzero
ordinary periods. For $1\le j\le k$, the spatial class $d=3^j$ contains
\begin{equation}\label{eq:three-count}
 2^{3^j-1}-2^{3^{j-1}-1}
\end{equation}
periodic words and has basin mass $2^{3^j}-2^{3^{j-1}}$.
\end{corollary}

\begin{proof}
The congruence $2^{3^{k-1}}\equiv-1\pmod{3^k}$ follows by induction:
cubing a number congruent to $-1$ modulo $3^k$ gives $-1$ modulo
$3^{k+1}$. Apply Corollary~\ref{cor:divisor-bound}. Only two terms in
\eqref{eq:Ad} survive for $d=3^j$, giving \eqref{eq:three-count}.
\end{proof}

The corollary provides a lower bound, not a claim of exactly $k$ periods
for every $k$. It separates $k$ classes and counts all their words without
determining the individual temporal periods within each class.

\section{The even-order cases of the conjecture}\label{sec:odd-order}

For an odd integer $u>1$, write $\ord_u(2)$ for the least positive
$r$ such that $2^r\equiv1\pmod u$.

\begin{lemma}[Constant-degree condition]\label{lem:degree}
If $\Spec(n)=\{T\}$, then
\begin{equation}\label{eq:constant-degree}
 \ord_d(2)=\ord_T(2)=\ord_n(2)
 \qquad\text{for every }d\mid n,\ d>1.
\end{equation}
\end{lemma}

\begin{proof}
Take an element $\alpha$ of exact multiplicative order $d$ in an
algebraic closure of $\F$. It is a root of $Q_n$, so
\eqref{eq:local-single} implies that $1+\alpha$ has exact order $T$.
An element of odd order $u>1$ has degree $\ord_u(2)$ over $\F$, since
this is the least exponent $r$ for which raising to $2^r$ fixes it.
But $\F(\alpha)=\F(1+\alpha)$, so the two degrees agree. Taking $d=n$
gives the final equality. All these orders are defined: $\alpha\ne0,1$,
so $1+\alpha\ne0,1$ and $T>1$ is odd.
\end{proof}

For composite $n$, condition \eqref{eq:constant-degree} is the
divisor-order characterization of an \emph{overpseudoprime to base $2$}
\cite[Theorem~2]{Overpseudoprimes}. Equal field degrees are necessary
for equal local periods, but they do not make the periods equal.

\begin{theorem}[Odd-order reduction]\label{thm:odd-order}
If an odd composite $n$ satisfies $|\Spec(n)|=1$, then $n$ is an
overpseudoprime to base $2$ and $\ord_n(2)$ is odd. Equivalently,
Conjecture~\ref{conj:single} holds for every odd $n$ for which
$\ord_n(2)$ is even.
\end{theorem}

\begin{proof}
The overpseudoprime condition follows from Lemma~\ref{lem:degree}.
Suppose the common order $r=\ord_n(2)$ were even. For each exact
prime-power divisor $p^e\parallel n$, that lemma gives
$\ord_{p^e}(2)=r$. The residue $2^{r/2}$ has order two modulo $p^e$.
For an odd prime power, the only nonidentity solution of $z^2=1$ is
$z=-1$: the two factors of $(z-1)(z+1)$ have no common odd prime
divisor, so one is divisible by $p^e$.
Thus $2^{r/2}\equiv-1\pmod{p^e}$ for every $p^e\parallel n$.
The Chinese remainder theorem gives $2^{r/2}\equiv-1\pmod n$.
Corollary~\ref{cor:divisor-bound} now supplies at least two nonzero
periods, a contradiction.
\end{proof}

\begin{remark}
Even $\ord_n(2)$ alone does not imply that $-1$ is a power of $2$
modulo $n$. For example, $\ord_{15}(2)=4$ but the powers of $2$
modulo $15$ omit $-1$. The proof uses the single-period hypothesis to
force equal orders on \emph{all} prime-power divisors before applying
the Chinese remainder theorem. Width $15$ already fails the degree
condition because $\ord_3(2)=2$ and $\ord_5(2)=4$.
\end{remark}

\section{Polynomial certificates and exact masses}\label{sec:certificates}

The remaining widths need not satisfy \eqref{eq:minus-one}. A fixed-space
test gives a second way to separate periods. For any odd $n>1$ and
$u\ge1$, define
\begin{equation}\label{eq:gcd-degree}
 g_n(u)=\deg\gcd\bigl(Q_n(x),(1+x)^u+1\bigr),
 \qquad F_n(u)=2^{g_n(u)}.
\end{equation}
All polynomial gcds are monic gcds over $\F$.

\begin{proposition}[A certificate of period diversity]\label{prop:certificate}
$F_n(u)$ is the number of periodic words whose ordinary period divides
$u$, including zero. There is more than one nonzero ordinary period
if and only if some $u$ satisfies
\begin{equation}\label{eq:proper-gcd}
 0<g_n(u)<n-1.
\end{equation}
For such a certificate, exactly
\begin{equation}\label{eq:certificate-counts}
 2^{g_n(u)}-1\quad\text{and}\quad 2^{n-1}-2^{g_n(u)}
\end{equation}
nonzero periodic words have, respectively, period dividing $u$ and
period not dividing $u$. The corresponding full basins contain twice
these numbers of words.
\end{proposition}

\begin{proof}
The fixed subspace of $D^u$ is the kernel of multiplication by
$(1+x)^u+1$ on $\F[x]/Q_n$. For polynomials $Q,f$, this kernel has
dimension $\deg\gcd(Q,f)$: writing $Q=gP$ and $f=gf_0$, with
$g=\gcd(Q,f)$, the congruence $fv=0\pmod Q$ is equivalent to
$P\mid v$, because $P$ and $f_0$ are coprime. This proves the count.

A proper nonzero fixed subspace contains a nonzero point fixed by $D^u$
and misses another, so those points cannot have equal periods. Conversely,
if $a<b$ are two occurring periods, choose $u=a$. A period-$a$ point is
fixed and a period-$b$ point is not, giving \eqref{eq:proper-gcd}.
Subtracting zero and the fixed-space size gives
\eqref{eq:certificate-counts}. The basin counts follow from
\eqref{eq:crt}.
\end{proof}

The fixed subspace is invariant under $R$ and $D$, and under their
inverses on $W_n$. It is therefore a union of complete joint orbits.
The certificate separates whole necklace components by the periods in
their cores; it does not merely select unrelated words. This is a direct
application of the partition structure. It needs no choice of a word in
either part and no irreducible factorization of $Q_n$. The exponentiation
in \eqref{eq:gcd-degree} can be reduced modulo $Q_n$ throughout.

\begin{theorem}[An explicit family of certificates]\label{thm:gcd-identity}
Assume \eqref{eq:minus-one}. For every positive integer $k$, put
$h=\gcd(n,k)$ and take $Q_1=1$. Then
\begin{equation}\label{eq:gcd-identity}
 \boxed{\gcd\bigl(Q_n(x),(1+x)^{kB}+1\bigr)=Q_h(x).}
\end{equation}
In particular, any divisor $1<d<n$ gives a certificate $u=dB$ with
$g_n(u)=d-1$.
\end{theorem}

\begin{proof}
Equation~\eqref{eq:ehrlich} means
$(1+x)^B=x^{-1}\pmod{Q_n}$. Since $x$ is a unit modulo $Q_n$,
the gcd in \eqref{eq:gcd-identity} equals $\gcd(Q_n,x^k+1)$.
The roots of $Q_n$ are the nonidentity $n$th roots of unity. Among
them, the roots of $x^k+1$ are precisely the nonidentity $h$th roots.
As $Q_n$ is squarefree, the monic gcd is $Q_h$.
\end{proof}

Thus the same arithmetic congruence yields both a separation theorem
and a fully specified algebraic certificate. For a proper divisor $d$,
the two groups of nonzero periodic words have sizes
\[
 2^{d-1}-1\quad\text{and}\quad2^{n-1}-2^{d-1}.
\]
Only the first population has temporal period dividing $dB$.

\subsection{Recovering exact periods}

Let $E_n(u)$ be the number of periodic words of exact ordinary period
$u$, with $E_n(1)=1$. M\"obius inversion gives
\begin{equation}\label{eq:exact-census}
 E_n(u)=\sum_{v\mid u}\mu(u/v)2^{g_n(v)}.
\end{equation}
There are $E_n(u)/u$ ordinary cycles of length $u$ and, for $u>1$,
$2E_n(u)$ initial words with eventual period $u$.
Under \eqref{eq:minus-one}, each of these ordinary cycles yields one
necklace component, with
\begin{equation}\label{eq:component-from-period}
 d=\gcd(n,u),\qquad t=u/d,\qquad
 \text{component mass}=2u.
\end{equation}
Without that hypothesis, the ordinary cycles can be joined by rotation,
and their number is not in general the number of necklace components.

For a known maximal period $P=\lcm_i\lambda_i$, a finite refinement of
the certificate criterion is
\begin{equation}\label{eq:prime-tests}
 |\Spec(n)|=1
 \quad\Longleftrightarrow\quad
 g_n(P/\ell)=0\ \text{for every prime }\ell\mid P.
\end{equation}
Indeed, every proper divisor of $P$ divides $P/\ell$ for some prime
$\ell\mid P$, while a point with full support attains period $P$.
This criterion assumes that $P$ is already known; it does not remove
the arithmetic difficulty of determining it.

\section{Two exact applications}\label{sec:examples}

\subsection{A spatial separation at width nine}

At $n=9$, take $s=3$, so $B=7$. The factorization is
\[
 Q_9=(x^2+x+1)(x^6+x^3+1).
\]
The second factor is irreducible because its roots have order $9$ and
$\ord_9(2)=6$. The orders of $1+x$ modulo the two factors are $3$ and
$63$, respectively. For the second, the order divides $2^6-1=63$;
the two prime-divisor tests are
\begin{align*}
 (1+x)^{21}&\equiv 1+x^3,&
 (1+x)^9&\equiv x^5+x^2+x
 \pmod{x^6+x^3+1}.
\end{align*}
Both residues differ from one, establishing exact order $63$.
Equation~\eqref{eq:lcm-spectrum} therefore gives
$\Spec(9)=\{3,63\}$.

\begin{table}[htbp]
\centering
\small
\renewcommand{\arraystretch}{1.3}
\begin{tabular}{@{}rrrrrr@{}}
\toprule
Spatial $d$ & Ordinary $\tau$ & Necklace $t$ & Core words & Components & Basin mass\\
\midrule
1 & 1 & 1 & 1 & 1 & 2\\
3 & 3 & 1 & 3 & 1 & 6\\
9 & 63 & 7 & 252 & 4 & 504\\
\bottomrule
\end{tabular}
\caption{The complete partition at width nine, grouped by ordinary period.
The first row is the zero core. Each component has one leaf per cycle
vertex. The basin masses sum to $512$.}
\label{tab:nine}
\end{table}

The zero component and the period-$3$ component have the same unlabeled
necklace graph: a loop with one leaf. Their masses differ, and only one
has a nonzero periodic core. Graph shape alone would lose precisely the
information needed for the conjecture.

Theorem~\ref{thm:gcd-identity} also gives the shorter certificate
\[
 \gcd\bigl(Q_9,(1+x)^{21}+1\bigr)=Q_3.
\]
Its degree is two, so it separates three nonzero core words from the
other $252$, without first determining both exact periods.

\subsection{An odd-order case with a complete census}\label{sec:mersenne}

The odd-order reduction leaves some genuine arithmetic work. For example,
$n=2047=23\cdot89$ satisfies the constant-degree condition with
$\ord_{2047}(2)=11$: both prime factors have order $11$. It is therefore
not excluded by Theorem~\ref{thm:odd-order}. The full-field structure of
Mersenne widths gives a separate exact result.

\begin{proposition}[Ordinary periods at Mersenne widths]\label{prop:mersenne}
Let $n=2^r-1>1$. For every positive $u$,
\begin{equation}\label{eq:mersenne-gcd}
 g_n(u)=\gcd(n,u)-1.
\end{equation}
Consequently
\begin{equation}\label{eq:mersenne-spectrum}
 \Spec(n)=\{d>1:d\mid n\},
 \qquad E_n(d)=\sum_{e\mid d}\mu(d/e)2^{e-1}
 \quad(d\mid n).
\end{equation}
In particular, Conjecture~\ref{conj:single} holds at every Mersenne width.
\end{proposition}

\begin{proof}
The roots of $Q_n$ are $\mathbb F_{2^r}\setminus\{0,1\}$.
Translation $\alpha\mapsto1+\alpha$ permutes this set. If
$h=\gcd(n,u)$, exactly $h$ elements of the cyclic group
$\mathbb F_{2^r}^\times$ satisfy $\beta^u=1$, and one is $\beta=1$.
There are therefore $h-1$ common roots in \eqref{eq:gcd-degree}.
Squarefreeness proves \eqref{eq:mersenne-gcd}.

The orders of the nonidentity elements of a cyclic group of order $n$
are exactly its divisors greater than one. These are the local orders
in \eqref{eq:lcm-spectrum}; taking least common multiples adds no other
values. This proves the spectrum. For $d\mid n$, substitute
$g_n(e)=e-1$ in \eqref{eq:exact-census}.
\end{proof}

For $d>1$ dividing a Mersenne width, the number of ordinary cycles is
\begin{equation}\label{eq:primitive-necklaces}
 \frac{E_n(d)}d=\frac1{2d}\sum_{e\mid d}\mu(d/e)2^e.
\end{equation}
It is half the number of binary necklaces of length $d$ with full spatial
period. The equality follows from the usual M\"obius count of primitive
words; it is a numerical correspondence, not an identification of the
two kinds of cycle.

At $n=2047$, the complete nonzero census is
\begin{equation}\label{eq:2047}
 \begin{array}{c|l}
 \text{Ordinary period}&\text{Number of periodic words}\\\hline
 23&2^{22}-1\\
 89&2^{88}-1\\
 2047&2^{2046}-2^{88}-2^{22}+1.
 \end{array}
\end{equation}
Together with the zero word these counts sum to $2^{2046}$; doubling
them accounts for all $2^{2047}$ starting words. Dividing each nonzero
entry by its period counts ordinary cycles. These counts need not equal
the numbers of necklace components; Theorem~\ref{thm:separation} does not apply here:
$\ord_{2047}(2)$ is odd, so no power of $2$ is $-1$ modulo $2047$.

The proportion of nonzero periodic words with period less than $2047$ is
exactly
\begin{equation}\label{eq:rare-periods}
 \frac{2^{88}+2^{22}-2}{2^{2046}-1}.
\end{equation}
Thus very small classes can decide the conjecture even when nearly all
core words have the maximal period. The exact partition detects both
exceptional classes without relying on their frequency. For example,
$u=23$ gives $g_{2047}(23)=22$, already a complete certificate of
period diversity.

\section{The remaining obstruction}\label{sec:remaining}

The results leave a precise problem. A composite counterexample must
satisfy the constant divisor-order condition
\eqref{eq:constant-degree}, with common order $r$ odd. This is a
necessary condition, not a sufficient condition and not an assertion
that any counterexample exists. Two further necessary bounds clarify
what an algebraic obstruction would have to rule out.

\begin{proposition}\label{prop:capacity}
If $\Spec(n)=\{T\}$ and $r=\ord_n(2)$, then
\begin{equation}\label{eq:capacity}
 n\mid T\mid 2^r-1,\qquad n-1\le\varphi(T),
\end{equation}
where $\varphi$ is Euler's totient function.
\end{proposition}

\begin{proof}
Every nonidentity $n$th root $\alpha$ belongs to $\mathbb F_{2^r}$,
and $1+\alpha$ has exact order $T$, giving $T\mid2^r-1$.
Choose $\alpha$ of exact order $n$. Both $1+\alpha$ and
$1+\alpha^{-1}$ have order $T$. Since
$1+\alpha^{-1}=\alpha^{-1}(1+\alpha)$, raising to $T$ gives
$\alpha^T=1$, hence $n\mid T$. Finally the $n-1$ distinct elements
$1+\alpha$, with $\alpha^n=1$ and $\alpha\ne1$, all have order $T$.
A cyclic group contains exactly $\varphi(T)$ elements of that order.
\end{proof}

In particular $n-1>\varphi(2^r-1)$ is another exclusion, since
$T\mid2^r-1$ implies $\varphi(T)\le\varphi(2^r-1)$. More generally,
the unresolved task can be expressed without listing any word:
\begin{quote}
For every odd composite $n$ satisfying the constant divisor-order
condition with odd common order, find a positive integer $u$ for which
$0<\deg\gcd(Q_n,(1+x)^u+1)<n-1$.
\end{quote}
By Proposition~\ref{prop:certificate}, solving this remaining task would
settle Conjecture~\ref{conj:single}. The hard part is a uniform argument
that forces such a proper fixed subspace; a formula that tests a chosen
$u$ is not that argument.

There is also inheritance across widths. If $d>1$ and $d\mid n$, repetition is
an injective map $V_d\to V_n$ commuting with $D$, and it preserves exact
ordinary periods, using the classical repetition construction
\cite[Theorem~5 and its proof]{Ehrlich}. Therefore
\begin{equation}\label{eq:inheritance}
 \Spec(d)\subseteq\Spec(n).
\end{equation}
An exclusion at one divisor excludes every multiple. To investigate
the remaining composites, one can first discard widths with an already
excluded proper divisor, then seek a certificate in the residual cases.
This inheritance does not assert that agreement at the prime divisors
forces agreement at their product: additional polynomial factors appear
at the larger width.

The application developed here is consequently a theorem and a reduction,
rather than a general solution of the conjecture. The spatial structure
of the necklace partition proves period diversity for an infinite
arithmetic class and, through the degree condition, excludes all
composite widths of even order. The fixed-space formulation supplies
exact certificates and population counts beyond that class. It isolates
the remaining difficulty as an interaction between addition and
multiplicative order in finite fields, while preserving a disjoint
accounting of every word of the power set.

\begin{thebibliography}{9}

\bibitem{Introduction}
\emph{Consequences of Binary Necklaces}, companion manuscript.

\bibitem{Partition}
\emph{The Partitioning of Power Sets Into Graphs}, companion manuscript.

\bibitem{BLM}
F.~Breuer, E.~L\"otter, and B.~van der Merwe,
\emph{Ducci-sequences and cyclotomic polynomials},
Finite Fields and Their Applications \textbf{13} (2007), no.~2,
293--304.
\href{https://doi.org/10.1016/j.ffa.2005.11.003}{doi:10.1016/j.ffa.2005.11.003}.

\bibitem{Ehrlich}
A.~Ehrlich,
\emph{Periods in Ducci's n-number game of differences},
The Fibonacci Quarterly \textbf{28} (1990), no.~4, 302--305.
\href{https://www.fq.math.ca/Scanned/28-4/ehrlich.pdf}{Full text}.

\bibitem{Overpseudoprimes}
V.~Shevelev, G.~Garc\'ia-Pulgar\'in, J.~M.~Vel\'asquez-Soto,
and J.~H.~Castillo,
\emph{Overpseudoprimes, and Mersenne and Fermat numbers as primover numbers},
Journal of Integer Sequences \textbf{15} (2012), Article~12.7.7.
\href{https://cs.uwaterloo.ca/journals/JIS/VOL15/Castillo/castillo2.html}
{Journal article}.

\bibitem{OEIS}
The OEIS Foundation Inc.,
\emph{The On-Line Encyclopedia of Integer Sequences}, entry A138004,
contributed by T.~D.~Noe.
\href{https://oeis.org/A138004}{oeis.org/A138004}.
The entry records the single-period primality conjecture; accessed
8 October 2026.

\end{thebibliography}
\end{document}
